\documentclass[14pt,a4paper]{extarticle}

% ---------- Paquetes ----------
\usepackage[utf8]{inputenc}
\usepackage[T1]{fontenc}
\usepackage[spanish]{babel}
\usepackage[margin=2.5cm]{geometry}
\usepackage{graphicx}
\usepackage{hyperref}
\usepackage{booktabs}
\usepackage{xcolor}
\newcommand{\completar}[1]{\textcolor{red}{[#1]}}

\hypersetup{colorlinks=true, linkcolor=black, urlcolor=blue}

\begin{document}

% ---------- Tapa ----------
\begin{titlepage}
\centering
% Para agregar el logo: subí la imagen a Overleaf y descomentá la línea de abajo
% \includegraphics[width=3cm]{logo_uba.png}\par\vspace{0.5cm}
{\LARGE\bfseries Universidad de Buenos Aires\par}
\vspace{0.2cm}
{\Large Facultad de Ciencias Económicas\par}

\vspace{0.2cm}
{\large\itshape Laboratorio de Métodos Cuantitativos Aplicados a la Gestión\par}
\vspace{0.5cm}

\begin{figure}[h]
  \centering
  \includegraphics[width=0.2\textwidth]{Logo-FCE.png}
\end{figure}

\vspace{0cm}
\rule{0.8\textwidth}{0.4pt}\par
\vspace{0.5cm}
{\huge\bfseries Influencia de la temperatura en la demanda eléctrica argentina\\[0.2cm] ¿Cuándo se vuelve crítica?\par}
\vspace{0.3cm}
\rule{0.8\textwidth}{0.4pt}\par
\vspace{0.6cm}
{\Large Trabajo Práctico Grupal\par}
\vspace{0.1cm}
{\large Exploración y análisis de datos asistido por IA\par}

\vspace{0.6cm}
\noindent\rule{0.7\textwidth}{0.4pt}


\vspace{0.6cm}
{\Large\scshape Integrantes\par}
\vspace{0.3cm}
{\large Berzoni Micaela · D'Angelo Genaro · Fiorini Catalina \\[0.15cm]
Mingorance Federico · Seri Rafael\par}

\vfill
\vspace{0.2cm}
{\normalsize Repositorio: \url{https://github.com/RafaelSeri/TP_LABO_MET}\par}
\vspace{0.2cm}
{\large 2do Cuatrimestre 2026\par}
\end{titlepage}

% ---------- Cuerpo (esqueleto) ----------
\section{Resumen}
Este trabajo analiza en qué medida la temperatura influye en la demanda de energía eléctrica en Argentina y en qué épocas del año esa influencia resulta más crítica. Se utilizan datos oficiales de demanda eléctrica y temperatura promedio entre enero de 2001 y agosto de 2026. El análisis, realizado en Python, incluye la limpieza de los datos, estadísticos agrupados por \completar{trimestre/estación} y visualizaciones de la evolución de la demanda y de su relación con la temperatura. Los resultados muestran que \completar{hallazgo principal}. Se concluye que \completar{conclusión principal}.

\section{Introducción}
La demanda de energía eléctrica en Argentina no se mantiene constante a lo largo del año. En los meses de verano, el uso masivo de equipos de aire acondicionado eleva el consumo, mientras que en invierno la calefacción eléctrica también puede generar aumentos. Estos picos ponen a prueba la capacidad del sistema eléctrico y pueden derivar en cortes de suministro, por lo que comprender sus causas resulta clave para la planificación energética. En este contexto, el presente trabajo busca responder la siguiente pregunta: \textit{¿en qué medida influye la temperatura en la demanda eléctrica argentina y en qué épocas del año esa influencia resulta más crítica?}

Para abordarla se utilizan datos de demanda eléctrica y temperatura elaborados por CAMMESA (Compañía Administradora del Mercado Mayorista Eléctrico) y publicados por el Ministerio de Economía de la Nación. \completar{A continuación, la Sección 3 describe los datos y la metodología empleada; la Sección 4 presenta los resultados; y la Sección 5 expone las conclusiones, limitaciones y posibles extensiones del trabajo.}


\section{Datos y metodología}

\subsection*{Fuente y selección de los datos}
Los datos se obtuvieron de la sección \textit{Economía en números}, que ofrece un archivo Excel de indicadores de nivel de actividad compuesto por 64 hojas. Se seleccionó únicamente la hoja 1.28, \textit{Demanda de Electricidad}, ya que reúne en un mismo cuadro la demanda eléctrica y la temperatura promedio.

\subsection*{Estructura y variables}
La hoja presenta la información organizada en tres bloques según su frecuencia: anual (2001--2025), trimestral (desde el primer trimestre de 2001 hasta el segundo de 2026) y mensual (desde enero de 2001 hasta agosto de 2026, con 308 registros). Para el análisis se utilizó el bloque \completar{mensual}. Las variables consideradas son:

\begin{itemize}
  \item \textbf{Demanda total} (GWh): energía eléctrica demandada en el período.
  \item \textbf{Demanda por tipo de usuario} (GWh): residencial, comercio e industria, y grandes usuarios. Esta apertura está disponible a partir de enero de 2005.
  \item \textbf{Temperatura promedio} (°C): temperatura media del período.
  \item \textbf{Potencia máxima} (MW): pico de potencia demandada en el período.
\end{itemize}

La hoja incluye además la demanda desagregada por empresa distribuidora y la potencia instalada por tipo de tecnología, que no fueron utilizadas por exceder el alcance de la pregunta.

\subsection*{Transformaciones}
Las series de demanda total y temperatura promedio no presentan valores faltantes. En cambio, la demanda por tipo de usuario carece de datos para los 48 meses comprendidos entre 2001 y 2004, por lo que los análisis que utilizan esta apertura se restringen al período 2005--2026. Sobre los datos se realizaron las siguientes transformaciones: \completar{qué limpieza hicieron: renombrar columnas, separar año y mes, eliminar filas vacías, etc.}. Para analizar la estacionalidad, los meses se agruparon en \completar{trimestres/estaciones del año}. Dado que 2026 cuenta solo con datos hasta agosto, \completar{se excluyó de las comparaciones anuales / se mantuvo aclarando que es un año incompleto}.

\subsection*{Herramientas}
El análisis se desarrolló en Python, utilizando Google Colab como entorno de trabajo y las librerías pandas para la manipulación de datos y matplotlib y seaborn para las visualizaciones. Se utilizaron herramientas de inteligencia artificial como apoyo en la escritura de código y la interpretación de resultados, cuyo uso se detalla en la Sección 6.

\section{Resultados}
% Gráficos exportados desde Colab, tablas e interpretación.
% Ejemplo para insertar un gráfico:
% \begin{figure}[h]
%   \centering
%   \includegraphics[width=0.8\textwidth]{grafico1.png}
%   \caption{Descripción del gráfico.}
%   \label{fig:grafico1}
% \end{figure}

\subsection*{Resumen estadístico}
\completar{Tabla con los estadísticos agrupados e interpretación.}

\subsection*{Evolución temporal de la demanda}

El análisis (Figura \ref{fig:evolucion}) nos permite observar el patrón de la demanda eléctrica a lo largo del año, comparándolo directamente con el ciclo de temperaturas.
La demanda no es lineal respecto a la temperatura. Se observan dos picos claros: uno en los meses más fríos (junio-julio) y otro en los meses más cálidos (diciembre-febrero). Por otro lado, el bajo consumo se da en los meses templados.

\begin{figure}[h]
  \centering
  \includegraphics[width= 1\textwidth]{demanda_relativa_promedio_mes.png} 
  \caption{Evolución de la demanda eléctrica diaria promedio y temperatura por mes.}
  \label{fig:evolucion}
\end{figure}

\subsection*{Comparación entre épocas del año}

Se observa en este gráfico de barras (Figura \ref{fig:estaciones}) la demanda relativa promedio por estación del año para los tres segmentos principales: Residencial, Comercio e Industria, y Grandes Usuarios.

El segmento Residencial muestra la mayor variación entre estaciones. Su demanda es máxima en invierno (por encima del 112\%) y mínima en primavera/otoño (por debajo del 90\%). Esto nos muestra que los hogares son los principales responsables de los picos estacionales del sistema.

El segmento de Comercio e Industria tiene su pico en verano (superando el 106\%), probablemente por la refrigeración en locales comerciales.

Los Grandes Usuarios presentan una demanda similar en todas las estaciones (cerca del 100\%), lo que demuestra que su consumo está atado a procesos productivos y no a las condiciones climáticas.

\begin{figure}[h]
  \centering
  \includegraphics[width= 1\textwidth]{demanda_relativa_estacion_segmento.png}
  \caption{Demanda relativa promedio por estación del año y segmento de usuario.}
  \label{fig:estaciones}
\end{figure}


\subsection*{Relación entre temperatura y demanda}

El gráfico de líneas (Figura \ref{fig:rangos}) muestra la demanda relativa promedio agrupada por rangos de temperatura (de menos de 12°C a más de 24°C).

Se observa claramente que la demanda mínima se concentra en el rango templado de 15 a 21°C (valores cercanos al 93\%). A medida que la temperatura se aleja, la demanda se dispara. El extremo frío (menos de 12°C) genera un pico de demanda total del 107.6\%, mientras que el extremo cálido (más de 24°C) genera un pico similar, de 107\%.

La línea naranja (Residencial) es la que tiene la pendiente más pronunciada en ambos extremos, confirmando que es el segmento que más reacciona al frío y al calor. La línea verde (Comercio e Industria) sube solo con el calor, y la roja (Grandes Usuarios) se mantiene casi plana.

\vspace{0.5cm} 

El gráfico de dispersión (Figura \ref{fig:dispersion}) muestra cada mes del dataset con un punto (Temperatura vs. Demanda Relativa), coloreado por estación, formando una parábola.

Se nota claramente cómo los puntos rojos (Verano) se agrupan a la derecha, lo cual habla de altas temperaturas y alta demanda; los azules (Invierno) a la izquierda, referido a las bajas temperaturas y alta demanda; y los verdes/naranjas (Primavera/Otoño) en el centro con demanda baja. La dispersión de los puntos en esta zona central (15-20°C) indican que, cuando el clima es templado, otros factores toman mayor relevancia relativa para las variaciones de demanda.

\begin{figure}[htbp]
  \centering
  \includegraphics[width= 1\textwidth]{demanda_temperatura.png}
  \caption{Demanda relativa según rango de temperatura.}
  \label{fig:rangos}
\end{figure}

\begin{figure}[htbp]
  \centering
  \includegraphics[width= 1\textwidth]{dispersion_temperatura.png}
  \caption{Relación temperatura-demanda según estación del año.}
  \label{fig:dispersion}
\end{figure}


\section{Conclusiones}
% Hallazgos principales, limitaciones y posibles extensiones.

\subsection*{Hallazgos principales}
\completar{Respuesta a la pregunta: cuánto influye la temperatura y en qué época es más crítica.}

\subsection*{Limitaciones}
\completar{2026 incompleto (hasta agosto); datos mensuales que ocultan olas de calor o frío de pocos días; temperatura promedio nacional; otros factores que afectan la demanda (actividad económica, tarifas, población).}

\subsection*{Posibles extensiones}
\completar{Qué se podría analizar a futuro.}

\section{Reflexión metodológica}
\completar{Uso de IA (herramientas, para qué, qué se corrigió), revisiones realizadas (intragrupal y externa) y aprendizajes.}

\begin{thebibliography}{9}
\bibitem{dataset}
CAMMESA. Demanda de Electricidad (Cuadro 1.28). Publicado por el Ministerio de Economía de la Nación, \textit{Economía en números}. \url{https://www.argentina.gob.ar/economia/datos/nivel-de-actividad}
\end{thebibliography}

\end{document}

